\documentclass[a4paper,12pt]{extarticle}
\usepackage[dvipsnames]{xcolor}
\usepackage{array,amsmath,fancybox,amsfonts,graphicx,amsthm,amssymb,wasysym,latexsym,slashed,anysize,subfigure}
\usepackage{tikz,calc}


\usepackage[spanish]{babel}
\usepackage[utf8]{inputenc}
\usepackage{fancyhdr,float,hyperref}
\usepackage{lettrine,multicol,pdflscape,enumerate}

%entorno para química: texto y math mode
%\newcommand*\chem[1]{\textup{#1}}
\newcommand*\chem[1]{\ensuremath{\mathrm{#1}}}

%formatos tachado
\usepackage{ulem}
%    \ sout {text}% tachado
%     \ uwave {text}% línea de onda
%     \ xout {text}% tachado diagonal
%     \ uuline {text}% Doble subrayado


%para los entornos de soluciones
\usepackage{mdframed}

	%uso: 	entorno (begin,end): {mdframed}[backgroundcolor=brown!10] 

%tamaños fuentes extra: 8,9,14,17,20
	\usepackage{extsizes}
	
 	%fuente librebaskerville
 	\usepackage[T1]{fontenc}
	\usepackage{librebaskerville}
	

	
	%fracciones $$ en texto más elaboradas
	 \usepackage{xfrac}
	 
	%permitir que una ecuación larga esté separada por un salto de página
	\allowdisplaybreaks[4]
		
	%gráficos varios 
	\usepackage{graphics}
	
	%fancy tables
	\usepackage{booktabs}
	\usepackage[small,centerlast,up]{caption}

	
	%mejora en la exposición de guarismos
	\usepackage{numprint}
	

\usepackage{everypage,draftwatermark} %Marca de agua
\SetWatermarkText{\hspace*{-1cm}\parbox{\linewidth}{\center B2\\dist}}
\SetWatermarkScale{2}
\SetWatermarkLightness{0.94}
\SetWatermarkAngle{45}



\DeclareMathOperator{\lin}{lin}
\DeclareMathOperator{\id}{\textbf{id}}
\DeclareMathOperator{\modulo}{mod}

\newcommand{\mc}[3]{\multicolumn{#1}{#2}{#3}}
\newcommand{\halmos}{\begin{flushright}
\rule{1mm}{2.5mm} \end{flushright}}
\renewcommand{\qed}{\halmos}


%versores: %uso: \uveci\ne\hat{\imath}_{\uvec{i}}+\uvec{j}+\uvec{k}$
\usepackage{bm}
\newcommand{\uveci}{{\bm{\hat{\textnormal{\bfseries\i}}}}}
\newcommand{\uvecj}{{\bm{\hat{\textnormal{\bfseries\j}}}}}
\DeclareRobustCommand{\uvec}[1]{{%
  \ifcsname uvec#1\endcsname
     \csname uvec#1\endcsname
   \else
    \bm{\hat{\mathbf{#1}}}%
   \fi
}}


%algún nuevo operador
\DeclareMathOperator{\cotan}{cotan}


%para textbf en entorno enumerate
\renewcommand{\labelenumi}{%
 \textbf{\theenumi}.
}

%anysize
\marginsize{.65cm}{.72cm}{.7cm}{.6cm}%requiere el paquete anysize,izq,dcha,arriba,abajo


\newtheorem{teo}{{\sc \textbf{Teorema}}}
\newtheorem{defi}{{\sc \textbf{Definición}}}
\newtheorem{coro}{{\sc \textbf{Corolario}}}


%opening
\title{\textbf{	}}
\date{}
%\pagestyle{fancy}
%\fancyhead[LO,RE]{
%					%Física, B2 
%%\href{http://tinyurl.com/fisicayquimica}{\nolinkurl{tinyurl.com/fisicayquimica}} 
%					}
%\fancyhead[RO,LE]{
%%Dept. Física y Química
%					}
%%\fancyhead[CO,CE]{Ricardo Torres\\http://tinyurl.com/iesCMGfyq}
%
\usepackage{enumitem}
\setlist[itemize]{leftmargin=3ex}
\setlist[itemize,1]{label=$\blacktriangleright$}
\setlist[itemize,2]{label=$\triangleright$}

\begin{document}
%\maketitle
\pagenumbering{gobble}

\vspace*{-1cm}

\begin{mdframed}[backgroundcolor=blue!10]  
\center
{\Large Temario}\\\normalsize{e hitos miliarios}
\end{mdframed}

\vspace{1.5cm}


\begin{multicols}{2}

	\begin{enumerate}
	\setcounter{enumi}{-1}
	
 		\subsection*{\fbox{Primer trimestre} 	}
	 	\item  {\small RESPUESTA ELÁSTICA}
	 	  \begin{itemize}
	 	    \item Oscilador armónico
	 	  \end{itemize}
	 	  
	 	\item {\small ONDAS ARMÓNICAS}
	 		
	 		\begin{itemize}
	 			\item Ondas transversales en una cuerda 
	 			\item Ondas de sonido 	
	 		\end{itemize}

	 	\item {\small ÓPTICA FÍSICA}
	 
	 			\begin{itemize}
	 				\item Caracterización y espectro {\scriptsize EM}
					\item Materiales transparentes
	 			\end{itemize}
		
			
	 	\item  {\small ÓPTICA GEOMÉTRICA}
		 		\begin{itemize}
		 			\item El proceso de transmisión-reflexión
	 				\item Óptica geométrica \textit{paraxial}  		
	 			\end{itemize}

	 	
	 \subsection*{\fbox{Segundo trimestre}}

	 	\item  {\small TEÓRIA CLÁSICA DE CAMPOS}
	 		\begin{itemize}
	 			\item Geometría y dinámica 			
				\item Campo gravitatorio y electrostático
				\item Electrodinámica
			\end{itemize}
	\vspace{-1ex}

	\begin{tikzpicture}
    \hspace*{-0.6cm}\draw [thick,dotted] (0,1) -- (9,1);
  \end{tikzpicture}
      
  
		 \item[\textbf{4\textonehalf}.] RELATIVIDAD ESPECIAL
      \begin{itemize}
        \item Principios de relatividad
		    \item Masa y energía
		  \end{itemize}

	\begin{tikzpicture}
    \hspace*{-0.6cm}\draw [thick,dotted] (0,1) -- (9,1);
  \end{tikzpicture}
   	\vspace{-4ex}
   
			
\subsection*{\fbox{Tercer trimestre}}

		\item {\small FÍSICA CUÁNTICA Y NUCLEAR}
		  \begin{itemize}
		    \item Dualidad de \textit{de Broglie}
		    \item Cinemática del efecto fotoeléctrico								
		    \item Desintegración radiactiva
		  \end{itemize}
						
 \end{enumerate}
\end{multicols}

\vspace{\fill}

\begin{center}
	\section*{Bibliografía}
\end{center}
	\begin{itemize}
%		\item Libro de texto propuesto.
%			\newline \textit{Util como referencia para la teoría y para los problemas.}
    \item Notas de clase.
    \newline {\color{RawSienna} Consulta y ampliación. La estructura es totalmente coherente con la de nuestro curso, aunque en ocasiones algunos puntos se tratan con mayor profundidad que en las clases.}
%		\item Alonso, Finn. {\footnotesize FÍSICA: CAMPOS Y ONDAS. VOL. II}. 
%			\newline {\color{RawSienna} Consulta y ampliación. El nivel del libro es aproximadamente el de un primer curso de grado, aunque asumible como consulta para este curso. Tomar \textit{cum grano salis}.}
	\end{itemize}

%\begin{center}
%	\section*{¿Qué es y qué no es {\large FÍSICA B2}?}
%\end{center}
%
%\begin{multicols}{2}
%	\fbox{ \textbf{ES}}
%
%		\begin{itemize}
%			\item \dots una oportunidad para saber más de un montón de fenómenos cotidianos.
%			\item \dots un laboratorio de tu curso de mates.
%			\item \dots para  \textit{llevar al día} y echarle horas.
%		\end{itemize}
%	%\vfill\null
%	\columnbreak
%%	\vfill\null		
%	
%	\begin{flushright}\fbox{ \textbf{NO ES}}\end{flushright}
%		\begin{itemize}	
%			\item \dots un curso \textit{riguroso}: estaremos  siempre en tensión con el instrumental matemático.
%			\item \dots una asignatura \textit{imposible} de aprobar. 
%			\item \dots una asignatura a \textit{dejar para el final}.
%		\end{itemize}
%\end{multicols}



%
%{\large
%
%\begin{center}
%	\section*{Sobre criterios y exámenes (\textit{extracto})}
%	\begin{itemize}
%		\item Cada trimestre: dos exámenes ($27\%$ y $63\%$); $10\%$ otros.
%		\begin{itemize}
%		\item[a)] La nota de final de curso se hará a partir del redondeo de la media de las notas \textit{no redondeadas} de las tres evaluaciones. Apruebas si es $\geq 5$.
%		
%		\item[b)] Si esta nota fuese $<5$ tendrás la oportunidad de recuperar enfrentando un \textit{examen final ordinario \textbf{global}}. \textit{Si apruebas} tu nota final se calcula entonces promediando la nota del examen final ordinario \textbf{con un cinco}.
%		
%%		\item[c)] Adicionalmente se establece la opción de que los alumnos que hayan aprobado y lo consideren oportuno puedan presentarse voluntariamente al examen final ordinario con objeto de aumentar su calificación. La nota en este caso \textbf{será la media aritmética} de la calificación obtenida en el apartado \textit{a)} y la nota obtenida en el examen final ordinario.
%		
%		\item[c)] Si suspendes en la ordinaria \textit{aúúún} puedes ir a la convocatoria extraordinaria (obviamente global). Tu nota final en este caso se calculará redondeando la nota final del examen extraordinario.
%		\end{itemize}
%	\end{itemize}
%\end{center}
%
%\vspace{1cm}
%
%\begin{center}
%	\section*{Sugerencias y comentarios}
%\end{center}
%\vspace{1ex}
%
%		\begin{itemize}
%
%			\item Si has llegado hasta aquí con una \textbf{nota inferior al notable} en {\small MATEMÁTICAS} o en {\small FÍSICA Y QUÍMICA} te aconsejo encarecidamente que dediques alguna tarde a reflexionar sobre el sentido de las expresiones \textit{tomar contacto con la realidad}, \textit{estudiante de bachillerato científico-técnico} y \textit{echarle más horas que tiempo tiene la orilla del mar.}
%			\item \textbf{No se propone libro de texto}: la asignatura tendrá teoría y práctica basada en la toma de apuntes, en la consulta de las \textit{notas del curso} y en el trabajo a partir de \textit{dossiers}.
%		 	\item \textbf{Recomiendo usar folios} y \textbf{organizar} todos y cada uno de los materiales que trabajemos: archivador/fundas transparentes/... Tú sabrás, que ya eres mayorcit@.
%		 	\item Aprender a \textbf{trabajar con autonomía} es uno de los grandes logros que debes alcanzar en bachillerato. Para ello, y aunque suene contradictorio, es \textit{esencial} \textbf{aprender a preguntar}, tanto en forma como en contenido: deja atrás cuanto antes la vergüenza, la timidez\dots y la galbana.
%		 	\item Como pasará en todas tus asignaturas el tiempo se nos echará encima: \textbf{aprovecha el inicio} de curso para preguntar y ponerte a nivel de trabajo. Te costará pero cuanto antes, mejor.
%%		 	\item Me ayudaría mucho que rellenarais un cuestionario inicial \textbf{anónimo} para saber algo más de vosotros y poder ajustar las clases al grupo. Llegáis a él a través de \href{https://tinyurl.com/incioB2}{https://tinyurl.com/incioB2}.
%		\end{itemize}
%	}

\end{document}

